% amsmath's matrix environments written as their macro pairs, which LaTeX
% accepts because \begin{bmatrix} runs \bmatrix: \bmatrix ... \endbmatrix,
% \Bmatrix, \vmatrix, \Vmatrix, and \smallmatrix. Each pair in a numbered
% equation, in \[ \], in inline math, in align*, with \\ rows in a numbered
% align and gather, through a macro of the document, and nested. Inline in a
% footnote, a theorem, and a heading with the table of contents, with
% hyperref, and with \\ rows in an eqnarray and an alignat.
\documentclass{article}
\usepackage{amsmath,amsthm}
\usepackage{hyperref}
\newtheorem{theorem}{Theorem}
\newcommand{\col}[2]{\Bmatrix #1 \\ #2 \endBmatrix}
\begin{document}
\tableofcontents
\section{Rotations by $\protect\bmatrix c & s \protect\endbmatrix$}

A row in a numbered equation:
\begin{equation}
v = \bmatrix 1 & 0 \endbmatrix
\end{equation}
The other pairs in a numbered equation:
\begin{equation}
p = \Bmatrix a & b \\ c & d \endBmatrix \vmatrix e & f \\ g & h \endvmatrix
  \Vmatrix i & j \\ k & l \endVmatrix \smallmatrix m & n \\ o & q \endsmallmatrix
\end{equation}
In a display:
\[ B = \bmatrix 1 & 2 \\ 3 & 4 \endbmatrix \Bmatrix 5 \\ 6 \endBmatrix
  \vmatrix 7 \endvmatrix \Vmatrix 8 \endVmatrix \smallmatrix 9 & 0 \endsmallmatrix \]
Inline, a column, $w = \bmatrix 1 \\ 0 \endbmatrix^T$, and the others,
$\Bmatrix r \\ s \endBmatrix$, $\vmatrix t & u \\ x & y \endvmatrix$,
$\Vmatrix z \\ \alpha \endVmatrix$, and $\left( \smallmatrix \beta & \gamma \\ \delta & \epsilon \endsmallmatrix \right)$.

In align*:
\begin{align*}
A &= \bmatrix a & b \\ c & d \endbmatrix & C &= \Bmatrix e \\ f \endBmatrix \\
D &= \vmatrix g & h \\ i & j \endvmatrix & E &= \Vmatrix k & l \\ m & n \endVmatrix \\
F &= \smallmatrix o & p \\ q & r \endsmallmatrix \\
P &= \Bmatrix e \\ f \endBmatrix & Q &= R \\
S &= \bmatrix a & b & c \endbmatrix & T &= U
\end{align*}
In a numbered align, with rows inside each matrix:
\begin{align}
G &= \bmatrix 1 \\ 2 \endbmatrix \label{eq:g} \\
H &= \Bmatrix 3 \\ 4 \endBmatrix \\
I &= \vmatrix a & b \\ c & d \endvmatrix \\
J &= \Vmatrix e & f \\ g & h \endVmatrix \\
K &= \smallmatrix 5 \\ 6 \endsmallmatrix
\end{align}
In a numbered gather:
\begin{gather}
L = \bmatrix a \\ b \endbmatrix \\
M = \Bmatrix c \\ d \endBmatrix \vmatrix e \\ f \endvmatrix \\
N = \Vmatrix 7 \\ 8 \endVmatrix \smallmatrix 9 \\ 0 \endsmallmatrix
\end{gather}
The align's first row is \eqref{eq:g}. Nested, in a numbered align:
\begin{align}
W &= \bmatrix \smallmatrix 1 \\ 2 \endsmallmatrix & 0 \\ 0 & 1 \endbmatrix \\
X &= 1
\end{align}
In an eqnarray and an alignat:
\begin{eqnarray}
Y &=& \bmatrix 1 \\ 2 \endbmatrix \\
Z &=& 3
\end{eqnarray}
\begin{alignat}{2}
a &= \Vmatrix 1 \\ 2 \endVmatrix &\quad b &= 4 \\
c &= 5 &\quad d &= 6
\end{alignat}
Through a macro, $u = \col{x}{y}$, and nested:
\[ O = \bmatrix \Vmatrix a \endVmatrix & \smallmatrix 1 & 2 \\ 3 & 4 \endsmallmatrix \\
  \vmatrix b \endvmatrix & \Bmatrix c \endBmatrix \endbmatrix \]
In a footnote.\footnote{The vector $\bmatrix 0 \\ 1 \endbmatrix$.}

\begin{theorem}
The matrix $\Vmatrix 2 & 3 \endVmatrix$ is a norm.
\end{theorem}

\section{Another section}
\end{document}
